\begin{thebibliography}{99}

\bibitem{Repo}
V. Reshetnikov, \emph{ProveIt: Analysis/Polylogarithms/docs},
source snapshot
\path{3a6d80ed618d839915deb0d19ab687f45e81d995}
(October 8, 2026).
\url{https://github.com/VladimirReshetnikov/ProveIt/tree/3a6d80ed618d839915deb0d19ab687f45e81d995/Analysis/Polylogarithms/docs}.
The accompanying source manifest identifies the individual files.

\bibitem{Panzer}
E. Panzer, \emph{The parity theorem for multiple polylogarithms},
Journal of Number Theory \textbf{172} (2017), 93--113.
\url{https://arxiv.org/abs/1512.04482}.
In particular Theorem 1.3, Corollary 1.4, and the depth-two formulas.

\bibitem{XuWang}
C. Xu and W. Wang, \emph{Dirichlet type extensions of Euler sums},
Comptes Rendus Math\'ematique \textbf{361} (2023), 979--1010.
\url{https://doi.org/10.5802/crmath.453}.
The specialization used here is equation (20), section 3.4,
printed page 994.

\bibitem{Kubert}
D. S. Kubert, \emph{The universal ordinary distribution},
Bulletin de la Soci\'et\'e Math\'ematique de France
\textbf{107} (1979), 179--202.
\url{https://doi.org/10.24033/bsmf.1891}.

\bibitem{Ouyang}
Y. Ouyang, \emph{On the universal norm distribution},
Journal of the Ramanujan Mathematical Society
\textbf{17} (2002), no.~4, 287--311.
\url{https://arxiv.org/abs/math/0210297}.

\bibitem{Coffey}
M. W. Coffey, \emph{Series representations for the Stieltjes constants},
Rocky Mountain Journal of Mathematics \textbf{44} (2014), 443--477.
Preprint: \url{https://arxiv.org/abs/0905.1111}.
The parameter-derivative formulas follow from its Hurwitz-zeta
addition and differentiation identities.

\bibitem{KowalskiNikeghbali}
E. Kowalski and A. Nikeghbali,
\emph{Mod-Poisson convergence in probability and number theory},
author manuscript.
\url{https://people.math.ethz.ch/~kowalski/mod-poisson.pdf}.
The permutation-cycle example explains the classical reciprocal-gamma
factor underlying the coefficient mechanism.

\bibitem{RomikScherer}
D. Romik and R. Scherer,
\emph{Alternative summation orders for the Eisenstein series \(G_2\)
and Weierstrass \(\wp\)-function},
arXiv:1811.01523, version 2 (2020).
\url{https://arxiv.org/abs/1811.01523}.

\bibitem{Ramanujan}
S. Ramanujan, \emph{On certain arithmetical functions},
Transactions of the Cambridge Philosophical Society
\textbf{22} (1916), no.~9, 159--184.
Transcription of the original paper:
\url{https://ramanujan.sirinudi.org/Volumes/published/ram18.html}.

\bibitem{ABP}
G. W. Anderson, W. D. Brownawell, and M. A. Papanikolas,
\emph{Determination of the algebraic relations among special
\(\Gamma\)-values in positive characteristic},
Annals of Mathematics \textbf{160} (2004), 237--313.
\url{https://annals.math.princeton.edu/wp-content/uploads/annals-v160-n1-p06.pdf}.
Section 1.4.2 distinguishes the classical Koblitz--Ogus sufficient
criterion from Rohrlich's conjectural converse; the paper's
positive-characteristic theorem is a different result.

\bibitem{DLMF}
NIST, \emph{Digital Library of Mathematical Functions},
sections 5.11 (gamma asymptotics), 23 (Weierstrass elliptic
functions), 25.11 (Hurwitz zeta), and 25.15 (Dirichlet \(L\)-functions).
\url{https://dlmf.nist.gov/}.
Consulted October 8, 2026.

\end{thebibliography}
